\section{Document Graph \texorpdfstring{\emojidoctitle}{Document Graph}}
\label{sec-doc}


A document graph typically models the connections between different documents or various granularity of documents. It is widely observed in real-world scenarios~\citep{schuhmacher2014knowledge, sonawane2014graph, xu2020document}, such as hyperlinks connecting different websites and citations linking one paper to another. Additionally, a document graph's relationships between sentences and entities can explicitly capture semantic and syntactic contextual information. The structural information in documents can serve as a valuable resource for GraphRAG, aiding LLMs in various tasks. For example, in the retrieval process in RAG, the document containing the answer may have less apparent connections with the question, such as in multi-hop question answering~\citep{mavi2024multi}. However, we can identify these relevant documents through their connections to other documents whose context is strongly aligned with the question's context~\citep{dong2024don}. In this section, we will systematically review the document graph.



\subsection{Application Tasks}
Document graphs are beneficial for a wide range of tasks. In this subsection, we will review the tasks where document graphs can play a significant role. Although document graphs have not yet been fully integrated with LLMs, they still offer great potential for enhancing the capabilities of LLMs in various applications.

\begin{itemize}[leftmargin=*]
\item \textbf{Multi-document Summarization(MDS)}: Multi-document summarization aims to condense the contents of multiple documents into a cohesive summary. Summarizing an entire corpus can involve large volumes of text, often exceeding the context window limitations of LLMs. Document graphs can help compress the corpus by extracting key components and their relationships, proving highly beneficial for MDS~\citep{yasunaga2017graph, wang2020heterogeneous, li2020leveraging, zhang2023contrastive, xie2022gretel, zhang2023enhancing, yang2018integrated, li2023compressed, chen2021sgsum}. These graphs also provide different levels of granularity for summarization through hierarchical clustering~\citep{edge2024local}.


\item \textbf{Text generation}: Text generation focuses on producing coherent and meaningful text. While text-based RAG models have been widely used to generate more reliable text, document graphs can further enhance this process by retrieving similar documents using graph topology. For instance, when writing the abstract of a paper, access to its cited papers in the related work section can significantly improve writing efficiency, as these referenced papers often contain relevant knowledge and context~\citep{wang2024augmenting, pan2020semantic}.
 
\item \textbf{Document Retrieval}: Document retrieval aims to find a list of documents relevant to a given query, which is a key task in Information Retrieval (IR). The exact query terms may not always appear together in the candidate document; however, by leveraging the connections between documents, we can retrieve related documents through those documents that are strongly linked to the query. Thus, it becomes essential to consider document-level relationships in the retrieval process. Several works~\citep{dong2024don, yu2021graph, zhang2018graph, liu2018matching} have leveraged document graphs with various granularities to improve document retrieval and ranking. Rather than retrieving the entire document, graphs can also be used to retrieve specific segments, such as chunks. 


\item \textbf{Document Classification:} Document classification is a fundamental task in natural language processing. Traditional methods often focus on the locality of words, limiting their ability to capture long-distance and non-consecutive word interactions. Moreover, these methods typically focus on individual documents, overlooking the relationships between them, where connected documents often exhibit homophily, meaning they are more likely to share similar labels. Building a graph can help enhance document classification by leveraging both local and global relationships between words within a single document~\citep{zhang2020every, liu2022hierarchical} and by utilizing document-level relationships~\citep{zhang2020text, xiao2022graph}.

\item \textbf{Question Answering}: Question answering aims to provide answers to questions based on information from documents and is a fundamental task for RAG. However, documents used for question answering can be long, and traditional methods often focus on local structures within these documents, neglecting their global structure, which is crucial for long-range understanding. Additionally, some multi-hop questions require reasoning across multiple documents, necessitating the use of document-level relationships.  Humans often consolidate scattered information into structured knowledge to streamline the reasoning process and make more accurate judgments, in line with cognitive load theory~\citep{sweller1988cognitive, li2024structrag}. Graph-based methods are well-suited for this task by constructing word-level document graphs~\citep{nie2022capturing, wang2023docgraphlm}, utilizing document-level relationships~\citep{he2024g, wang2024knowledge} and leveraging hierarchical interactions~\citep{fang2019hierarchical}.


\item \textbf{Relation Extraction:} Relation extraction aims to extract semantic relationships between entities in text, which often requires local, global, syntactic, and semantic dependencies, especially in the case of document-level relation extraction. Graphs have been proven to be helpful for document-level relation extraction~\citep{wang2020global, nan2020reasoning, sahu2019inter, christopoulou2019connecting, zhou2020global} by capturing these dependencies more effectively.

In addition to the aforementioned tasks, graphs have also been shown to be helpful in other areas, such as fake news detection~\citep{hu2021compare}, coherence assessment~\citep{liu2023modeling}, and machine translation~\citep{xu2020document}.
\end{itemize}


\subsection{Document Graph Construction}

Different tasks may require different types of document graphs, such as document-level, sentence-level, or word-level graphs. Therefore, the method of obtaining the document graph plays a crucial role. There are primarily two approaches to constructing a document graph: explicit construction and implicit construction.

\begin{itemize}[leftmargin=*]
\item \textbf{Explicit Construction.}  In real-world scenarios, many documents have explicit connections. Examples include web pages linked through hyperlinks, academic papers citing other works, and social media posts connected by reposts, comments, and interactions.  In these cases, it is natural to connect the documents to build a document graph, as the connected documents often share semantic relationships~\citep{peng2024connecting}. For instance, \citet{asai2019learning} constructed a Wikipedia graph based on hyperlinks between Wikipedia articles. \citet{li2022dynamic, li2021graph} follow the same way to leverage hyperlinks to construct graphs. \citet{yu2021kg} built a graph using an external knowledge graph (KG), where each node represents a retrieved passage mapped to entities within the KG, and two passage nodes are connected if their mapped entities are linked in the KG. 
These connections have been leveraged to pretrain large language models (LLMs), enhancing their performance across various tasks~\citep{yasunaga2022linkbert, zou2023pretraining, hu2021relation}.


\item \textbf{Implicit Construction.} Despite the presence of explicit connections between documents, different components within those documents also exhibit important semantic and syntactic relationships. The structure of these components can be highly beneficial for various tasks in natural language processing~\citep{wu2023graph, mihalcea2011graph}. For example, semantic parsing graphs, such as Abstract Meaning Representation (AMR) graphs, can be leveraged to enhance information extraction~\citep{zhang2021abstract, huang2017zero} and text summarization~\citep{dohare2017text}. Additionally, syntactic parsing trees are exploited to understand grammatical structure and resolve ambiguities~\citep{zhang2023survey}.

The implicit construction of document graphs is highly diverse, adapting to the specific requirements of different tasks. For instance, nodes in document graphs can represent various granularity, such as words, entities, sentences, text segments, paragraphs, documents, or topics. In addition, document graphs often exhibit heterogeneity, where edges connect different types of nodes. Next, we will detail the construction of different types of edges:

\begin{itemize}
    \item {\it Word-word edge:} The word-word edges connect words with semantic or syntactic relations or dependencies. There are various methods to establish these connections, including word co-occurrence within a sliding window~\citep{wang2020heterogeneous, yu2021graph, zhang2020text, liu2018matching, de2018question}, dependency parsing graphs~\citep{xu2018graph2seq, tang2020integration, xu2020document, zhang2019aspect, qiu2019dynamically, tu2020select} generated by NLP parsing tools~\citep{lee2011stanford, kumawat2015pos}, Abstract Meaning Representation~\citep{xu2021dynamic, zhang2021abstract, huang2017zero, dohare2017text, xu2021exploiting} and semantic graphs~\citep{dong2024don, song2018graph, pan2020semantic} based on different representations~\citep{banarescu2013abstract, shi2019simple}. Additional techniques include coreference resolution~\citep{xu2020document, sahu2019inter, leskovec2005extracting, dhingra2018neural, song2018exploring}, word embedding similarity~\citep{wang2023docgraphlm, wang2022me} and using large language models (LLMs) to extract relationships between words~\citep{edge2024local}. For entities, the construction process is similar to that for words, though entity extraction methods are often required to first identify entities~\citep{nasar2021named, al2020named, pawar2017relation}.

    \item {\it Word-Sentence edge:} Word-sentence edges connect words to sentences based primarily on the relationship of belonging. These edges are constructed by linking words to the sentences they appear in~\citep{wang2020heterogeneous, hu2021compare, christopoulou2019connecting}, with edge weights often measured using term frequency-inverse document frequency (TF-IDF). Besides, \citet{ramesh2023single} connects the entities with passage titles via hyperlinks with existing out-of-box entity linkers.

    \item {\it Sentence-Sentence edge:} Sentence-sentence edges connect sentences based on their semantic similarity or relationships. For example, these edges can be constructed through sentence interactions~\citep{yasunaga2017graph, maheshwari2024presentations, huang2021breadth}, similarities between TF-IDF representations~\citep{li2020leveraging, chen2021sgsum}, BM25~\citep{min2019knowledge}, sentence embeddings~\citep{liu2023modeling, wang2023docgraphlm, zhang2024coarse}, Part of speech (Pos) feature and N-Gram feature~\citep{yang2018integrated}. Additionally, \citet{zheng2020srlgrn} construct a Semantic Role Labeling(SRL) graph using  AllenNLP-SRL model~\citep{shi2019simple}. This allows for connecting sentences within long documents or across different documents, which is particularly useful for LLMs with limited context length.

    \item {\it Sentence-document edge:} Sentence-document edges are constructed by linking sentences to the specific documents they belong to~\citep{thayaparan2019identifying}.

    \item {\it Document-document edge:} Document-document edges connect documents based on their similarities. These edges can be constructed when entities in one document are referenced or shared by another~\citep{thayaparan2019identifying, dong2024don}, through document clustering~\citep{wang2023docgraphlm}, embedding similarity~\citep{li2020leveraging}, topic similarity~\citep{zhang2024coarse} or structural similarity~\citep{liu2023modeling}.
\end{itemize}
\end{itemize}

Document graphs typically exhibit heterogeneity, consisting of multiple types of edges. In addition, several hierarchical graphs have been proposed~\citep{weninger2012document, zhang2024teleclass}, where different levels of abstraction are captured. The graph structure can also be dynamic or updated during the learning process~\citep{wang2020heterogeneous, nan2020reasoning, wang2023docgraphlm, liu2022hierarchical}. Moreover, the node in the document graph can also be the response of LLMs. For example, GoR~\citep{zhang2024graph} connects the history responses of LLMs with the responding chunks of documents for long-context summarization.




\subsection{Retriever}

The retrievers for document graphs typically follow the general retriever design, as described in Section~\ref{sec:retriver}. However, there are some special retrieval methods as follows:

\begin{itemize}[leftmargin=*]
    \item \textbf{Pre-Retrieval:} As introduced earlier, there are many scenarios where building a document graph is necessary. However, constructing a fine-grained graph for a large volume of documents can be inefficient and unnecessary. The pre-retrieval aims to first retrieve relevant documents based on the query and then construct a graph. For example, \citet{thayaparan2019identifying} use pre-trained GloVe vectors to first extract relevant sentences and then construct a graph based on the retrieved sentences. \citet{zheng2020srlgrn, yu2021kg} also construct graphs based on the retrieved information.
    
    \item \textbf{Graph similarity-based retriever:}  Graph similarity aims to measure the similarity between two graphs. If both the query and the retrieval data source are graphs, calculating the graph similarity is essential to retrieve relevant information. For instance, \cite{zhang2018graph} leverages the General Maximum Common Subgraph (GMCS) method to retrieve related graphs.
    
    \item \textbf{Iterative Retriever:} In certain tasks, such as multi-step question answering, the node containing the answer might not be directly similar to the query. Iterative retrievers address this by first retrieving nodes related to the query and then using the retrieved information to iteratively retrieve subsequent nodes. For example, \citet{wang2024knowledge, zhang2021answering} utilize iterative retrieval methods for multi-step question-answering tasks, and \citet{ma2024think} iteratively retrieve the documents based on the existing knowledge graph. \citet{asai2019learning} train a Recurrent Neural Network (RNN) to recurrently retrieve relevant information with the Bayesian Personalized Ranking (BRR) loss.
    
    \item \textbf{Topology-based Retriever:} Various topological relationships in a graph can be used to measure different types of similarity. For example, proximity-based topological similarity measures the structural distance between two nodes, while role-based topological similarity assesses the similarity of the roles nodes play within the graph. These similarities can also be leveraged in the retrieval process~\citep{wang2024augmenting}.
\end{itemize}


\subsection{Organizer}

In this section, we describe the organizers for document graphs. GraphRAG on document graphs is still in its early stage, and many works do not incorporate explicit organizers, as shown in Section~\ref{sec:organizer}. We summarize the existing methods below:

\begin{itemize}[leftmargin=*]
    \item \textbf{Graph Pruning}: Graph Pruning refines the retrieved subgraph to reduce irrelevant information and improve computational efficiency. For example, \citet{hemmati2023multi} prune graphs based on the local clustering coefficient, while \citet{zhang2018graph1} employs path-centric pruning to incorporate off-path information. \citet{li2022graph} dynamically drop irrelevant nodes during decoding, and \citet{angelova2006graph} prune edges based on a similarity threshold. Additionally, \citet{edge2024local} uses community detection algorithms to create distinct communities that are then fed into the generators.
    \item \textbf{Reranking}:  Reranking methods aim to reorder retrieved information to facilitate the generation. For example, \citet{yu2021kg}, \citet{zhang2021answering}, and \citet{dong2024don} use GNNs to rerank retrieved passages. \citet{li2021graph} perform listwise reranking to reorder passages expanded via the graph structure.
\end{itemize}

\subsection{Generator}
In this section, we summarize commonly used generators for document graphs. Various methods are applied within document graphs, depending on the input format. Some approaches use the entire graph as input, in which case GNNs and Graph Transformers are commonly employed. Other methods take individual sentences as input, where RNNs or (Large) Language Models are suitable. Additionally, some works leverage both graph and text as inputs, requiring integrated methods that can process multimodal data effectively. We detail each category in the following:


\begin{itemize}[leftmargin=*]
    \item \textbf{GNN-based generator}: Numerous works model tasks as graph-related problems, leveraging GNNs as generators. Conventional GNNs, such as GCN~\citep{kipf2016semi}, GraphSAGE~\citep{hamilton2017inductive}, and GAT~\citep{velivckovic2017graph}, are widely adopted in various studies~\citep{peng2024connecting, thayaparan2019identifying, yasunaga2017graph, munikoti2023atlantic, pan2020semantic, dong2024don}. When graphs contain edge relations, Relational Graph Convolutional Networks (R-GCNs) are often used to capture these relationships effectively~\citep{de2018question, min2019knowledge}. In addition, some works incorporate graph contrastive learning techniques to enhance performance~\citep{xie2022gretel, zhang2022ke, hemmati2023multi}.
    \item \textbf{Graph Transformer-based generator}: Graph transformers~\citep{yun2019graph}, which capture global information across the graph, are used to encode graph structures for various tasks~\citep{zhang2020text, mei2021graph}. These models leverage the transformer’s self-attention mechanism to capture dependencies beyond local neighborhoods, making them well-suited for tasks requiring global context.
    \item  \textbf{RNN-based generator}: Recurrent Neural Networks (RNNs), such as LSTMs, are popular for processing sequences. Therefore, RNNs are employed when the input is in text form~\citep{song2018exploring}.
    \item \textbf{LLM-based generator}: LLM-based generators typically transform the retrieved subgraph into text before using large language models (LLMs) for generation~\citep{wang2024knowledge, edge2024local}. Various models are used in this approach. For instance, BERT~\citep{devlin2018bert} has been applied by \cite{tu2020select, li2022mutually, zhang2021answering} to support generation tasks, while \citet{yu2021kg} and \citet{ju2022grape} leverage the T5 model~\citep{raffel2020exploring}, and \citet{chen2020improving} use RoBERTa~\citep{liu2019roberta}
    \item \textbf{Integrated Generator}: Some works leverage both graph and text data simultaneously for generation, using integrated generators that combine graph models with text generation models. For example, \citet{xu2021dynamic} and \citet{wang2023docgraphlm} employ a combination of RoBERTa and GCN, while \citet{ramesh2023single} fuse GNN with T5 to harness the strengths of both graph structure and language model capabilities.
\end{itemize}

In addition to the aforementioned generators, some approaches embed graphs directly into text generation models. For example, \citet{li2020leveraging} proposes replacing traditional self-attention layers in transformers with graph-informed self-attention, enabling the model to integrate graph structure directly into the generation process. 








\subsection{Resources and Tools}
The data resources for document GraphRAG can include any type of document, and thus, we do not provide a comprehensive list here. Instead, in the following, we introduce several tools specifically designed or commonly used for GraphRAG on document graphs:



\begin{itemize}[leftmargin=*]
    \item \textbf{CoreNLP: }\footnote{\url{https://github.com/stanfordnlp/CoreNLP}} The Stanford CoreNLP natural language processing toolkit offers a comprehensive set of natural language processing tools, including token and sentence boundaries, parts of speech, named entities, numeric and time values, dependency and constituency parses, coreference, sentiment, quote attributions, and relations. These tools are valuable for constructing document graphs. 
    \item \textbf{spaCy: }\footnote{\url{https://github.com/explosion/spaCy}} The spaCy is an advanced natural language processing library known for its speed and neural network models, which are optimized for tasks such as tagging, parsing, named entity recognition, text classification, part-of-speech tagging, dependency parsing, sentence segmentation, lemmatization, morphological analysis, and entity linking. 
    \item \textbf{BLINK}: \footnote{\url{https://github.com/facebookresearch/BLINK}} BLINK is an entity linking python library that uses Wikipedia as the target knowledge base.
    \item \textbf{OpenIE}\footnote{\url{https://github.com/dair-iitd/OpenIE-standalone}}: Open Information Extraction (OpenIE) is a tool for extracting structured information from text. It is valuable for generating triples from unstructured text, which can be used as nodes and edges in document graphs.
    \item \textbf{CogComp NLP}\footnote{\url{https://github.com/CogComp/cogcomp-nlp}}: Developed by the Cognitive Computation Group, this suite includes tools for natural language processing tasks such as named entity recognition, sentiment analysis, and coreference resolution.
    \item \textbf{GraphRAG}~\citep{edge2024local}\footnote{\url{https://github.com/microsoft/graphrag}}: GraphRAG is a data pipeline and transformation suite that is designed to extract meaningful, structured data from unstructured text using the power of LLMs. The process involves extracting a knowledge graph out of raw text, building a community hierarchy, generating summaries for these communities, and then leveraging these structures when performing RAG-based tasks.
    \item \textbf{LangChain}~\footnote{\url{https://www.langchain.com/}}: LangChain is a framework for developing applications powered by LLMs, with use cases in document analysis and summarization, RAG, chatbots, and code analysis. Notably, LangChain supports Graph Transformers, which convert documents into graph-structured formats, making it highly suitable for processing document graphs.
    \item \textbf{Neo4j}\footnote{\url{https://neo4j.com/}}: Neo4j is a graph database platform offering comprehensive tools for storing, visualizing, managing, and querying graph data. It includes an LLM Graph Builder, which can extract graphs with LLMs. It also provides GraphRAG demos to demonstrate how to implement a LLMs and RAG system with Neo4j.
    \item \textbf{LlamaIndex}~\citep{Liu_LlamaIndex_2022}\footnote{\url{https://www.llamaindex.ai/}}: LlamaIndex is a data framework designed to support the development of applications based on LLMs. It allows developers to seamlessly integrate data sources, ranging from various file formats to applications and databases, with LLMs. LlamaIndex features an efficient data retrieval and query interface, enabling developers to input any LLM prompt and receive context-rich, knowledge-enhanced outputs. Notably, it includes a Property Graph Index that facilitates the building, modeling, storage, and querying of graphs.
    \item \textbf{Haystack}\footnote{\url{https://haystack.deepset.ai/}}: Haystack is an end-to-end framework for building applications powered by LLMs, Transformer models, vector search, and more. It supports a variety of use cases, including RAG, document search, question answering, and answer generation. Haystack enables the orchestration of state-of-the-art embedding models and LLMs into custom pipelines for creating comprehensive NLP applications. Additionally, it supports Neo4j as a DocumentStore, making it suitable for document graph storage and query operations.
\end{itemize}